\textbf{Proceso de desarrollo:} El proyecto se desarrolló siguiendo un enfoque iterativo incremental. Se inició con la especificación de requerimientos (documento SRS) que delineó las funcionalidades y objetivos clave. A partir de allí se realizó un diseño de arquitectura y se planificó el desarrollo en módulos: primero el backend y la estructura de datos, luego la app móvil de conductores y la web de usuarios, y finalmente el panel administrativo. Se emplearon controles de versión (GitHub) y trabajo colaborativo entre programadores. Durante la implementación, se llevaron a cabo pruebas unitarias en las funciones críticas del backend (por ejemplo, servicios de autenticación y envío de mensajes) y pruebas de integración para verificar la correcta comunicación entre la app móvil y el servidor en distintos escenarios (conectividad presente, pérdida/reconexión de red, etc.). Dado el énfasis en \textbf{tiempo real}, se probó específicamente el mecanismo de actualización continua: se simuló el movimiento de un vehículo usando coordenadas de prueba y se confirmó que los clientes recibían las actualizaciones sin retraso apreciable. Para las pruebas de rendimiento, se configuró un entorno de staging donde múltiples clientes suscribieron a los tópicos MQTT de ubicaciones; se midió el consumo de recursos del servidor y se verificó que el sistema mantuviera tiempos de respuesta bajos con hasta 10 clientes simultáneos (limitados por el alcance de la prueba). Esto brinda una indicación inicial de escalabilidad; otros estudios han reportado que soluciones similares soportan cargas mayores, e.g. un módulo basado en MQTT manejó 100 dispositivos concurrentes sin fallos de comunicación en pruebas controladas\href{https://fliphtml5.com/lqsof/cdpg/Art\%C3\%ADculos/\#:~:text=permitiendo\%20un\%20flujo\%20de\%20datos,del\%20sistema\%20fue\%20un\%20rigurosa}{{[}1{]}}, por lo que se espera que UrbanTracker pueda escalar a cientos de usuarios simultáneos con una infraestructura adecuada.

En resumen, la metodología combinó el desarrollo de software basado en requerimientos definidos y las pruebas enfocadas en los criterios de aceptación clave (actualización en ≤10 s, latencia ≤2 s, usabilidad, etc.). A continuación, se detallan los resultados obtenidos, incluyendo métricas de desempeño del sistema y el grado de cumplimiento de los objetivos planteados.
